%============================================================
% CH2 · LESSON 2 — COLOURS, ENTRY AND EXIT TIMES
%============================================================
\subsection{Colours, Entry and Exit Times}
\label{ssec:ch2-timestamps}

The \texttt{visited} array answers one question: \emph{has DFS been here?} For everything else in this chapter we need a finer answer: \emph{has DFS been here, and is it still here?} The call stack of Lesson~\ref{ssec:ch2-call-stack} suggests three states, and a clock that ticks whenever DFS enters or leaves a vertex turns them into numbers.

\subsubsection*{Three colours}

\begin{definitionbox}[White, gray, black]
During a DFS every vertex has one of three colours:
\begin{center}
\legendchip{gv}{\textbf{white} --- not entered yet} \qquad \legendchip{gvseen}{\textbf{gray} --- entered, not left (on the stack)} \qquad \legendchip{gvdone}{\textbf{black} --- left (finished)}
\end{center}
A vertex turns gray when \texttt{dfs(u)} starts and black when it returns. It never changes colour again. The gray vertices, at any moment, are exactly the vertices of the call stack --- a path from the root to the current vertex.
\end{definitionbox}

In code the colours are one \texttt{int} per vertex, \texttt{0/1/2}; the old \texttt{visited[u]} is simply ``colour $\ne$ white''. The difference between gray and black is what Lesson~\ref{ssec:ch2-cycles} needs to detect cycles in directed graphs.

\subsubsection*{Entry and exit times}

Keep a global counter \texttt{timer}. When DFS enters $u$, record $\tin{u} = \texttt{timer++}$; when it leaves $u$, record $\tout{u} = \texttt{timer++}$. With one clock for both events, a DFS of $n$ vertices uses the times $0, 1, \dots, 2n-1$, each exactly once.

\begin{lstlisting}[style=snippet]
void dfs(int u) {
    color[u] = 1; tin[u] = timer++;          // u turns gray
    for (int v : adj[u])
        if (color[v] == 0) { parent[v] = u; dfs(v); }
    color[u] = 2; tout[u] = timer++;         // u turns black
}
\end{lstlisting}

\begin{diagrambox}
\begin{tikzpicture}[x=0.62cm, y=0.62cm]
    \tikzset{gv/.append style={minimum size=6mm, font=\scriptsize\bfseries}}
    \begin{scope}[g2-v1/.style=gvseen, g2-v2/.style=gvseen, g2-v3/.style=gvdone, g2-v4/.style=gvcur,
                  g2-e12/.style={gtree, -{Stealth[length=2.2mm]}}, g2-e23/.style={gtree, -{Stealth[length=2.2mm]}},
                  g2-e24/.style={gtree, -{Stealth[length=2.2mm]}}]
        \pic {g2-D};
        \node[font=\scriptsize, anchor=west] at (-0.5,-2.6) {\textbf{time 4:} enter 4 \quad stack \texttt{1 2 4}};
    \end{scope}
    \begin{scope}[xshift=4.6cm, g2-v1/.style=gvseen, g2-v2/.style=gvdone, g2-v3/.style=gvdone, g2-v4/.style=gvdone,
                  g2-v5/.style=gvseen, g2-v6/.style=gvseen, g2-v7/.style=gvcur,
                  g2-e12/.style={gtree, -{Stealth[length=2.2mm]}}, g2-e23/.style={gtree, -{Stealth[length=2.2mm]}},
                  g2-e24/.style={gtree, -{Stealth[length=2.2mm]}}, g2-e15/.style={gtree, -{Stealth[length=2.2mm]}},
                  g2-e56/.style={gtree, -{Stealth[length=2.2mm]}}, g2-e67/.style={gtree, -{Stealth[length=2.2mm]}}]
        \pic {g2-D};
        \node[font=\scriptsize, anchor=west] at (-0.5,-2.6) {\textbf{time 9:} enter 7 \quad stack \texttt{1 5 6 7}};
    \end{scope}
    \begin{scope}[xshift=9.2cm, g2-v1/.style=gvdone, g2-v2/.style=gvdone, g2-v3/.style=gvdone, g2-v4/.style=gvdone,
                  g2-v5/.style=gvdone, g2-v6/.style=gvdone, g2-v7/.style=gvdone, g2-v8/.style=gvcur,
                  g2-e12/.style={gtree, -{Stealth[length=2.2mm]}}, g2-e23/.style={gtree, -{Stealth[length=2.2mm]}},
                  g2-e24/.style={gtree, -{Stealth[length=2.2mm]}}, g2-e15/.style={gtree, -{Stealth[length=2.2mm]}},
                  g2-e56/.style={gtree, -{Stealth[length=2.2mm]}}, g2-e67/.style={gtree, -{Stealth[length=2.2mm]}}]
        \pic {g2-D};
        \node[font=\scriptsize, anchor=west] at (-0.5,-2.6) {\textbf{time 14:} enter 8 \quad stack \texttt{8}};
    \end{scope}
\end{tikzpicture}
\figcap{Three moments of the DFS on $D$. Orange $=$ current vertex, light orange $=$ gray (on the stack), blue $=$ black. Thick blue arrows are the edges along which DFS entered a new vertex.}
\end{diagrambox}

Running the DFS on $D$ (outer loop over $s = 1..8$) gives:
\begin{center}
\small
\begin{tabular}{@{}l*{8}{c}@{}}
\toprule
vertex $u$ & 1 & 2 & 3 & 4 & 5 & 6 & 7 & 8 \\
\midrule
$\tin{u}$ & 0 & 1 & 2 & 4 & 7 & 8 & 9 & 14 \\
$\tout{u}$ & 13 & 6 & 3 & 5 & 12 & 11 & 10 & 15 \\
parent & --- & 1 & 2 & 2 & 1 & 5 & 6 & --- \\
\bottomrule
\end{tabular}
\end{center}

\begin{diagrambox}
\begin{tikzpicture}[x=0.78cm, y=0.5cm]
    \draw[->, darkgray] (-0.3,0) -- (16.2,0) node[right, font=\scriptsize] {time};
    \foreach \t in {0,...,15} {\draw[darkgray] (\t,0.12) -- (\t,-0.12); \node[font=\tiny, below] at (\t,-0.12) {\t};}
    % intervals: vertex/tin/tout/level/colour
    \foreach \v/\a/\b/\l/\c in {1/0/13/5/primaryblue, 2/1/6/4/successgreen, 3/2/3/3/accentorange, 4/4/5/3/accentorange,
                                5/7/12/4/successgreen, 6/8/11/3/accentorange, 7/9/10/2/softpurple, 8/14/15/5/primaryblue}{
        \draw[\c, line width=2.2pt] (\a,\l*0.62) -- (\b,\l*0.62);
        \draw[\c, line width=1pt] (\a,\l*0.62-0.25) -- (\a,\l*0.62+0.25);
        \draw[\c, line width=1pt] (\b,\l*0.62-0.25) -- (\b,\l*0.62+0.25);
        \node[font=\scriptsize\bfseries, fill=white, inner sep=1pt, text=\c] at ({(\a+\b)/2},\l*0.62+0.42) {\v};
    }
    \node[font=\ttfamily\small, anchor=west] at (-0.3,-1.5) {( 1 ( 2 ( 3 3 ) ( 4 4 ) 2 ) ( 5 ( 6 ( 7 7 ) 6 ) 5 ) 1 ) ( 8 8 )};
\end{tikzpicture}
\figcap{The intervals $[\tin{u}, \tout{u}]$ on the time line. Writing ``\texttt{(}$u$'' at $\tin{u}$ and ``$u$\texttt{)}'' at $\tout{u}$ gives a correctly nested bracket sequence --- hence the name of the next theorem.}
\end{diagrambox}

\subsubsection*{The parenthesis theorem}

\begin{theorembox}[Parenthesis theorem]
For any two vertices $u \ne v$ reached by a DFS, exactly one of the following holds:
\begin{enumerate}[leftmargin=*]
    \item the intervals $[\tin{u}, \tout{u}]$ and $[\tin{v}, \tout{v}]$ are \textbf{disjoint}, and neither vertex is a descendant of the other in the DFS forest;
    \item $[\tin{v}, \tout{v}] \subset [\tin{u}, \tout{u}]$, and $v$ is a \textbf{descendant} of $u$;
    \item $[\tin{u}, \tout{u}] \subset [\tin{v}, \tout{v}]$, and $u$ is a descendant of $v$.
\end{enumerate}
Two intervals never overlap partially.
\end{theorembox}

\textbf{Proof.} Say $\tin{u} < \tin{v}$ (otherwise swap the names). Look at the moment $\tin{v}$ when $v$ is entered.
\emph{Case $u$ is still gray then}, i.e.\ $\tin{v} < \tout{u}$. Gray vertices form the current stack, which is a path from the root through $u$ down to the vertex that is entering $v$ --- so $v$ is entered while \texttt{dfs(u)} is still open and becomes a descendant of $u$. A call returns only after all calls it made have returned, so \texttt{dfs(v)} returns before \texttt{dfs(u)}: $\tout{v} < \tout{u}$. That is case 2.
\emph{Case $u$ is already black}, i.e.\ $\tout{u} < \tin{v}$: the intervals are disjoint. $v$ cannot be a descendant of $u$ (descendants are entered while $u$ is gray), and $u$ cannot be a descendant of $v$ (it was entered first). That is case 1. \hfill$\blacksquare$

\begin{insightbox}[The \bigO{1} ancestor test]
\[
    u \text{ is an ancestor of } v \text{ (or } u = v) \iff \tin{u} \le \tin{v} \;\text{ and }\; \tout{v} \le \tout{u}.
\]
One DFS (\VE) and then every query ``is $u$ above $v$ in the DFS tree?'' costs two comparisons. On $D$: $1$ is an ancestor of $7$ ($0 \le 9$, $10 \le 13$); $2$ is not an ancestor of $6$ ($\tout{6} = 11 > 6$).
\end{insightbox}

\begin{warningbox}[Ancestor in the DFS tree is not the same as ``can reach'']
In $D$ there is an edge $5 \to 4$, so 5 can reach 4 --- but 4 is \emph{not} a descendant of 5: its interval $[4,5]$ lies before $[7,12]$. The DFS reached 4 earlier, from 2. The parenthesis theorem is about the tree that this particular DFS built, not about reachability. (For a rooted \emph{tree} given as input, the two notions coincide, which is why the test is so useful in Patterns~2 and~3.)
\end{warningbox}

\subsubsection*{The white-path theorem}

Which vertices become descendants of $u$? Not ``all vertices reachable from $u$'' (vertex 4 above), but almost:

\begin{theorembox}[White-path theorem]
$v$ becomes a descendant of $u$ in the DFS forest if and only if, at time $\tin{u}$, there is a path from $u$ to $v$ that consists of \textbf{white} vertices only.
\end{theorembox}

\textbf{Proof sketch.} ($\Rightarrow$) The tree path from $u$ to $v$ consists of vertices entered after $u$, so they were white at $\tin{u}$. ($\Leftarrow$) Suppose some vertex of the white path is not a descendant of $u$; let $w$ be the first such vertex and $x$ its predecessor on the path ($x = u$ or a descendant of $u$). DFS examines the edge $x \to w$ before $x$ is left; at that moment $w$ is either still white --- then it becomes a child of $x$ --- or it was entered after $\tin{u}$ and before $\tout{x} \le \tout{u}$, which by the parenthesis theorem makes it a descendant of $u$. Either way a contradiction. \hfill$\blacksquare$

On $D$ at time $\tin{5} = 7$, vertex 4 is already black, so the path $5 \to 4$ is not white and 4 is not a descendant of 5; the path $5 \to 6 \to 7$ is white, and indeed 6 and 7 end up below 5. The white-path theorem is the key step in the correctness proofs of cycle detection (Lesson~\ref{ssec:ch2-cycles}) and, much later, of strongly connected components (Chapter~19).

\subsubsection*{A second convention: pre-order numbers and subtree sizes}

The ``one clock'' times above are perfect for comparisons. For \emph{storing} subtree data in an array, a second convention is more convenient: number the vertices $0, 1, \dots, n-1$ in pre-order (call it $\mathrm{pos}[u]$) and compute the subtree sizes $\mathrm{sz}[u]$. Then the subtree of $u$ occupies exactly the block
\[
    \mathrm{pos}[u],\ \mathrm{pos}[u] + 1,\ \dots,\ \mathrm{pos}[u] + \mathrm{sz}[u] - 1
\]
of the pre-order, and $u$ is an ancestor of $v$ iff $\mathrm{pos}[u] \le \mathrm{pos}[v] < \mathrm{pos}[u] + \mathrm{sz}[u]$.

\begin{diagrambox}
\begin{tikzpicture}[x=1cm, y=1cm]
    \foreach \v/\i in {1/0, 2/1, 3/2, 4/3, 5/4, 6/5, 7/6, 8/7}{
        \node[gcell, minimum size=8mm, font=\small\bfseries] (c\i) at (\i*0.9,0) {\v};
        \node[font=\tiny, text=darkgray] at (\i*0.9,-0.6) {pos \i};
    }
    \draw[primaryblue, line width=1.4pt, decorate, decoration={brace, amplitude=5pt}] (-0.4,0.5) -- (5.8,0.5)
        node[midway, above=5pt, font=\scriptsize, text=primaryblue] {subtree of 1: pos $0..6$ (sz 7)};
    \draw[successgreen, line width=1.2pt, decorate, decoration={brace, amplitude=4pt, mirror}] (0.5,-0.85) -- (3.1,-0.85)
        node[midway, below=4pt, font=\scriptsize, text=successgreen] {subtree of 2: $1..3$};
    \draw[successgreen, line width=1.2pt, decorate, decoration={brace, amplitude=4pt, mirror}] (3.2,-0.85) -- (5.8,-0.85)
        node[midway, below=4pt, font=\scriptsize, text=successgreen] {subtree of 5: $4..6$};
\end{tikzpicture}
\figcap{The pre-order of $D$ is $1,2,3,4,5,6,7,8$ and the sizes are $(7,3,1,1,3,2,1,1)$. Every subtree is a contiguous block. Turning subtree questions into range questions on this array (``flattening'') is the subject of Chapter~13 (tree fundamentals); answering range queries on it fast is Chapter~15 (subtree and path queries). Here we only need the ancestor test and the block itself (Problem~\ref{ssec:ch2-military}).}
\end{diagrambox}

\subsubsection*{Code: timestamps and ancestor queries}

\begin{lstlisting}[style=cpp]
#include <bits/stdc++.h>
using namespace std;

// Colours, entry/exit times (one clock) and O(1) ancestor queries.
// Input: n m, m directed edges, then q queries "u v": is u an ancestor of v in the DFS forest?
vector<vector<int>> adj;
vector<int> color, tin, tout, par;
int timer = 0;

void dfs(int u) {
    color[u] = 1; tin[u] = timer++;              // white -> gray
    for (int v : adj[u])
        if (color[v] == 0) { par[v] = u; dfs(v); }
    color[u] = 2; tout[u] = timer++;             // gray -> black
}

bool isAncestor(int u, int v) {                  // u == v counts as an ancestor
    return tin[u] <= tin[v] && tout[v] <= tout[u];
}

int main() {
    int n, m;
    cin >> n >> m;
    adj.assign(n + 1, vector<int>());
    for (int i = 0; i < m; i++) {
        int a, b;
        cin >> a >> b;
        adj[a].push_back(b);
    }
    color.assign(n + 1, 0); tin.assign(n + 1, -1); tout.assign(n + 1, -1); par.assign(n + 1, 0);
    for (int s = 1; s <= n; s++)
        if (color[s] == 0) dfs(s);
    int q;
    cin >> q;
    while (q--) {
        int u, v;
        cin >> u >> v;
        cout << u << " above " << v << "? " << (isAncestor(u, v) ? "yes" : "no") << "\n";
    }
}
\end{lstlisting}

\begin{lstlisting}[style=trace]
input : D, then 4 queries   1 7   2 6   5 4   8 7
output: 1 above 7? yes
        2 above 6? no
        5 above 4? no        (5 -> 4 is an edge, but 4 was entered from 2)
        8 above 7? no        (8 -> 7 is an edge, but 7 was finished long before)
\end{lstlisting}

\subsubsection*{Check your understanding}

\begin{enumerate}[leftmargin=*]
    \item How many different values does \texttt{timer} take during a DFS of $n$ vertices? What are $\tin{r}$ and $\tout{r}$ for the root $r$ of a DFS from a single start that reaches all $n$ vertices?
    \item For a leaf $u$ of the DFS tree, what is $\tout{u} - \tin{u}$? For a general vertex, express $\tout{u} - \tin{u}$ through the subtree size.
    \item Can $\tin{u} < \tin{v} < \tout{u} < \tout{v}$ ever happen? Why not?
    \item In $D$, list every vertex that is gray at time 10.
    \item With the pre-order convention, which pre-order positions does the subtree of 6 occupy in $D$?
\end{enumerate}

\begin{solutionbox}[Answers]
\begin{enumerate}[leftmargin=*]
    \item $2n$ values, $0..2n-1$: every vertex is entered once and left once. The root is entered first and left last: $\tin{r} = 0$, $\tout{r} = 2n-1$.
    \item A leaf is left right after it is entered: $\tout{u} - \tin{u} = 1$. In general all $\mathrm{sz}[u]$ vertices of the subtree are entered and left inside $u$'s interval, which uses $2\,\mathrm{sz}[u]$ ticks: $\tout{u} - \tin{u} = 2\,\mathrm{sz}[u] - 1$. (On $D$: vertex 5 has $12 - 7 = 5 = 2 \cdot 3 - 1$.)
    \item No --- that would be a partial overlap, excluded by the parenthesis theorem: if $v$ is entered while $u$ is gray, $v$ is left before $u$.
    \item At time 10 vertex 7 is being left; just before that tick the stack was \texttt{1 5 6 7}. After it, gray are \textbf{1, 5, 6} (7 has just turned black).
    \item $\mathrm{pos}[6] = 5$ and $\mathrm{sz}[6] = 2$: positions \textbf{5 and 6} (vertices 6, 7).
\end{enumerate}
\end{solutionbox}
